% ECONOMICS_PROBLEM: {"id": "D82-280", "title": "Inequality-aware allocation", "jel_code": "D82", "source_id": "OP-0296"}
% !TeX program = xelatex
\documentclass[11pt,a4paper]{article}
\usepackage[margin=23mm]{geometry}
\usepackage{fontspec}
\setmainfont{Latin Modern Roman}
\usepackage{amsmath,amssymb,mathtools}
\usepackage{xcolor,enumitem,booktabs,longtable,array,xurl,hyperref}
\definecolor{muted}{HTML}{53636C}
\hypersetup{colorlinks=true,urlcolor=blue,linkcolor=blue}
\setlength{\parindent}{0pt}
\setlength{\parskip}{5pt}
\newcommand{\R}{\mathbb R}
\newcommand{\E}{\mathbb E}
\newcommand{\N}{\mathbb N}
\newcommand{\ind}{\mathbf 1}
\newcommand{\Var}{\operatorname{Var}}
\newcommand{\Cov}{\operatorname{Cov}}
\newcommand{\supp}{\operatorname{supp}}
\DeclareMathOperator*{\argmin}{arg\,min}
\DeclareMathOperator*{\argmax}{arg\,max}
\newcommand{\entrymeta}[3]{{\sffamily\small\color{muted}\textbf{\texttt{#1}}\quad|\quad #2\par #3\par}}
\newcommand{\Problemref}[1]{\texttt{#1}}
\newcommand{\JELref}[2]{\texttt{#2} (JEL \texttt{#1})}
\begin{document}
\section*{Inequality-aware allocation}
\textbf{Problem ID:} \texttt{D82-280}\quad \textbf{JEL:} \texttt{D82}\par
% BEGIN ECONOMICS STATEMENT
\label{problem:OP-0296}
{\sffamily\small\color{muted}Original subject group: Mechanism design and market design\par}
\label{definitions:problem:30007578}
\textbf{Audit classification:} Published research agenda. Dworczak’s verified §4 items 2–4 explicitly raises simple mechanisms, moments/robustness and multidimensional extensions. The displayed minimax problem is editorial.
\subsection*{Definitions and precise research target}
\textbf{Model and research target.} A unit mass of agents demand at most one homogeneous indivisible good. Type \((r,w)\) comprises private willingness to pay \(r\in[0,1]\) and social welfare weight \(w\in[0,W]\), where finite \(W>0\) is known. The designer has supply \(q\in(0,1)\) and knows only specified moments of the joint type distribution; let \(\mathcal F\) be all distributions consistent with them. A direct mechanism specifies allocation probability \(x(r)\in[0,1]\) and payment \(p(r)\), with utility \(r x(r)-p(r)\), truthful reporting, voluntary participation, and feasibility under every \(F\in\mathcal F\). With shadow value \(\lambda\geq0\) for public revenue, evaluate
\[
 V_F(M)=E_F[w(r x(r)-p(r))+\lambda p(r)],\qquad
 \sup_M\inf_{F\in\mathcal F}V_F(M).
\]
Characterize robust optimal mechanisms and when price menus, subsidies, or rationing improve this redistributive objective. Derive sharp guarantees for posted-price mechanisms relative to the full robust optimum, under clearly stated budget constraints and welfare-weight bounds. Explain why privately varying \(w\) affects the designer's objective without being directly reportable in utility; additional screening dimensions require an extended model.

\emph{Definitions and maintained assumptions (editorial unless a source definition is named).} Use a unit-mass atomless population with type law \(F\) on \([0,1]\times[0,W]\). The allocation depends only on reported willingness to pay \(r\); \(w\) is a designer welfare weight, not a private utility argument being truthfully screened. DSIC is \(rx(r)-p(r)\ge rx(r')-p(r')\) for all \(r,r'\), and IR is \(rx(r)-p(r)\ge0\). Impose \(-\bar s\le p(r)\le\bar p\) and a declared balanced-budget or public subsidy constraint; without a subsidy limit, high welfare weights can make unlimited payments to agents artificially optimal. Supply means \(\mathbb E_F[x(r)]\le q\) for every \(F\in\mathcal F\), interpreted as aggregate fractions in the atomless model. Specify which bounded moment functions and exact moment values define nonempty \(\mathcal F\). The objective treats a unit of public revenue at shadow value \(\lambda\) and agents’ net surplus at weight \(w\); these are normative primitives. Approximation ratios need a nonnegative normalization and a strictly positive optimum; otherwise use additive losses.
\subsection*{Variants and assumption changes}
\begin{enumerate}
\item \textbf{Moment ambiguity (source agenda, editorial model).} Fix means and covariance of \((r,w)\), bounded transfers and a subsidy budget. Characterize the minimax DSIC menu and whether a posted price achieves a guaranteed fraction of its positive value.
\item \textbf{Multiple qualities (source agenda, editorial model).} Replace the homogeneous good by finitely many qualities with explicit utility \(v(q,\theta)\) and reportable multidimensional type \(\theta\). Study robust mechanisms with redistributive weights; the one-dimensional monotonicity characterization no longer suffices.
\end{enumerate}
\subsection*{References and source audit}
\begin{enumerate}
\item §4 items2–4, printed pp.90–91. Supports the stated research area; exact displayed specification is editorial. \url{https://www.sigecom.hosting.acm.org/exchanges/volume_22/1/issue.pdf}
\end{enumerate}
\begin{enumerate}\raggedright
\item\hypertarget{definitions:source:30007578:1}{} Piotr Dworczak. 2024. \emph{Inequality and Market Design}. §4 items2–4, printed pp.90–91.
\url{https://www.sigecom.hosting.acm.org/exchanges/volume_22/1/issue.pdf}.
\end{enumerate}

\subsection*{Common conventions: Source mathematical conventions}

These conventions apply to each entry unless explicitly replaced.

\textbf{Models and identification.} A primitive model \(m\) specifies all probability laws, feasible choices, information, equilibrium and measurement rules used by its target. The admissible class \(\mathcal M\) is fixed independently of outcomes. If \(P_O(m)\) is its observable probability law and \(\tau(m)\) the target, its identified set at observable law \(P\) is
\[
 \mathcal I_\tau(P)=\{\tau(m):m\in\mathcal M,\ P_O(m)=P\}.
\]
Point identification means that this set is a singleton. An empty set signals incompatibility of data and assumptions. Coordinate bounds need not describe the joint set. Bounds are sharp only if every retained target is attainable or arbitrarily approachable by compatible models. Finite bounds require explicit restrictions on unobserved quantities; unrestricted confounding, hidden wealth, extrapolation or arbitrary equilibrium selection can yield uninformative sets.

\textbf{Causal interventions.} A potential outcome \(Y_i(p)\) is the outcome under a complete policy regime \(p\), including timing, financing, information and market spillovers. The target population and units are fixed before treatment. With interference, use the complete assignment vector or a justified exposure mapping; individual no-interference notation alone cannot identify an economy-wide intervention. A difference of mean potential outcomes does not require the cross-world joint distribution. Conditioning on post-treatment migration, survival, enrollment or employment changes the estimand and needs separate justification.

\textbf{Statistical statements.} An estimator, confidence set, forecast or test specifies a sampling law, model class and loss/coverage criterion. Uniform \((1-\alpha)\)-coverage means \(\inf_{m\in\mathcal M}\Pr_m\{\tau(m)\in C_n\}\geq1-\alpha\), or an explicitly asymptotic version. The whole parameter space trivially has coverage, so informativeness requires a stated width, risk or power objective. A forecasting improvement is relative to a named benchmark, information set and evaluation population.

\textbf{Equilibrium and optimization.} An equilibrium comprises feasible plans, optimality, beliefs where needed, budget constraints and all market/resource clearing. The primitives or a stated benchmark define each correspondence \(\mathcal E(p;m)\). Nonemptiness is assumed or proved, never inferred just from compact policy sets. A selected-equilibrium target uses a declared selection rule; a robust target quantifies over all permitted equilibria. Use suprema or infima unless attainment is established. An \(\varepsilon\)-optimal policy is feasible and within \(\varepsilon\) of the finite optimum value. Report an empty feasible set separately.

\textbf{Welfare and accounting.} Normative utility, interpersonal normalization, social weights, numeraire, discounting and the use of public receipts are primitives. Behavioral utility need not coincide with the welfare criterion. Household welfare includes ownership income and rents once; adding profits again to compensating variation generally double-counts them. A monetary equivalent is defined by an explicit transfer and price convention, with monotonicity and range conditions. Average effects, mechanism contributions, transfers and welfare losses are different targets. Infinite sums require convergence and dynamic feasible plans require terminal or no-Ponzi conditions.

\textbf{Source versions.} A working-paper date, revision date and journal publication year are distinguished. A DOI for a journal version is not automatically the DOI of an earlier manuscript. Locators refer to the stated version and printed pages where specified. A metadata-only check does not verify a claim about a full-text conclusion. Every entry's source subsection narrows the provenance claim accordingly.
% END ECONOMICS STATEMENT
\end{document}
